\section{Experiment Plan and Status}\label{app:plan}

Table~\ref{tab:plan} lists every experiment of the protocol (\datafile{experiments/experiments.yaml}), whether it ran,
and where its data are. Hours and commands of the full and minimal plans are in \datafile{experiments/PLAN.md}. The
experiments that did not run keep their workloads, runners and data folders in the repository, and running one makes
its table appear in the paper on the next build.

\paragraph{The campaign.}
The campaign that ran is a reduced version of the protocol, called the \emph{lean plan} in the repository
(\datafile{experiments/runs/lean.sh}). Its first session runs on one 2x RTX PRO 6000 pod of at most
\dv{lean.hard_cap_minutes} minutes, paying boot, build and the weight download once. It covers the setup, E2 reduced to
the two EXL3 boots without offload, E8 with scenarios A and B at all seven concurrency levels, two repetitions and the
probes, E10 on fixed-seed subsets of GSM8K, IFEval and HumanEval, and E3 reduced to the \texttt{gen30} plan against
upstream's zero-copy offload at the \texttt{gen30} host bytes, with one repetition at $N\in\{1,8,32\}$. A final
optional step adds a larger nested plan, the union of the \texttt{gen30} and \texttt{agent30} cold sets. The runner
checks the remaining time before every step, so that a session ends cleanly instead of in the middle of a
measurement.

\paragraph{How each session went.}
The first session ran on \dv{lean.session.date} for \dv{lean.session.minutes} minutes
(\datafile{data/controlled/LEAN/session.json}). It completed every step, including the optional one, plus a repeat of
E10 after the continuation fix of Section~\ref{sec:quality}. Two bootstrap bugs forced two restarts before the first
valid boot: git refused the copied sources, and a patch check could not tell whether overlapping patches had been
applied. Both are fixed in \datafile{stack/build/bootstrap_pod.sh}. The second session, on \dv{lean.session2.date}
(\dv{lean.session2.minutes} minutes, \datafile{experiments/runs/long.sh}), ran on a different host. A first attempt
was aborted by a third bootstrap bug, a git configuration step that ran before git was installed, and its boots were
invalid and stayed on the pod. The session then ran the long-context experiment E12, the E3 \texttt{gen30} point again
as a baseline on the same pod, and the zero-copy rung with the KV block size it needs. The third session
(\dv{lean.session3.minutes} minutes, \datafile{experiments/runs/agent.sh}) ran on a third host, whose record notes the
first host's public IP but a different CPU model and GPU bus ids (Table~\ref{tab:pod-hardware}). It ran the
coding-agent experiment E13 after a first attempt that stopped before any measurement because the image lacked git.

\paragraph{The quality reference.}
The quality reference is the official DeepSeek API, whose \texttt{deepseek-flash} model is served by
DeepSeek-V4.1-Flash according to its pricing page (fetched 2026-10-06, \datafile{data/static/pricing/}). The run sent
\dv{ctl.E11.api.requests} requests, all answered, through the proxy that logs the real token counts.

\begin{table}[h]
\centering\small
\begin{tabular}{llp{6.2cm}p{3.6cm}}
\toprule
id & GPUs & question & status \\
\midrule
E0 & none & static byte accounting & done (Section~\ref{sec:system}) \\
E1 & CPU & reproducible build of every variant, image digests and provenance & not run \\
E2 & 2x PRO 6000 & feasibility without host offload, exact failure & ran (session 1, EXL3 rungs only) \\
E3 & 2x PRO 6000 & host share, upstream zero-copy offload at equal host bytes & ran (two shares, sessions 1 and 2, zero-copy cannot serve) \\
E4 & 2x PRO 6000 & placement policy at 30\%: recipe plans, random, our profile plan & not run \\
E5 & 2x PRO 6000 & Engram int4 vs FP8 in host memory & not run \\
E6 & 2x PRO 6000 & SM budget of the host-expert stream & not run \\
E7 & 2x PRO 6000 & DSpark on vs off, code vs prose & not run \\
E8 & 2x PRO 6000 & main benchmark: $N$ = 1 to 64, probes, memory, PCIe & ran (session 1, closed loop, two repetitions in B, one in A beyond $N=1$ and in the probes) \\
E9 & 2x B300 & upstream reference with everything in GPU memory, original checkpoint & not run \\
E10 & 2x PRO 6000 & quality of the EXL3 stack & ran (session 1, fixed subsets) \\
E11 & 2x B300 & quality of the reference & ran through DeepSeek's API instead \\
E12 & 2x PRO 6000 & long context at 128K and 256K tokens & ran (session 2) \\
E13 & 2x PRO 6000 & coding-agent workload, streams at 40 tokens/s & ran (session 3) \\
\bottomrule
\end{tabular}
\caption{Controlled experiments and their status. Sessions: \datafile{data/controlled/LEAN/session.json},
\datafile{data/controlled/LEAN/session2.json}, \datafile{data/controlled/LEAN/session3.json}.}
\label{tab:plan}
\end{table}
